\chapter{Conclusion}

This dissertation benchmarked symmetric cryptography inside a zkVM, compared AES
and LowMC behavior under proof-oriented cost models, and implemented a final
authenticated AES pipeline for verifiable execution.

The work combines three layers that are often studied separately: measured
proving performance, implementation-level engineering choices, and security
reasoning for authenticated transcripts. Treating these layers together was
necessary because practical deployment decisions in zkVM environments depend on
their interaction, not on any one layer in isolation.

The main empirical findings are:

\begin{itemize}
\item in the report benchmark run,
\texttt{lowmc-r0-optimised} is about \(34.85\times\) slower than
\texttt{aes-r0} on proving time (517.209 s vs 14.843 s);
\item LowMC optimisation improved local performance materially
(\(\sim 25.59\%\) proving-time reduction versus baseline), but did not reverse
the final ranking against AES;
\item Adding HMAC increases user-cycle work at every tested block size
(about \(+1.38\%\) to \(+50.15\%\) depending on payload), while total-time
deltas are noisier in single-trial values.
\end{itemize}

The final repository structure reflects these findings. CTR-only and CTR+HMAC
paths are represented as separate benchmark target families
(\texttt{aes-ctr} and \texttt{aes-ctr-hmac}), enabling direct and auditable
mode-overhead comparisons in the same proving environment.

Methodologically, the dissertation highlights a practical lesson for zkVM
engineering: optimisation effort should be guided by trace-backed measurement and
artifact discipline, not by algorithmic intuition alone.

The project therefore contributes both a concrete implementation result and a
repeatable evaluation workflow. With commit-locked multi-trial reruns and
cross-zkVM replication, this baseline can be extended into stronger general
claims for authenticated cryptographic computation in verifiable-execution
systems.

\section{Objective-by-Objective Evaluation}

Table~\ref{tab:conclusion-objective-evaluation} gives the final status of each
measurable objective against the evidence presented in the report.

\begin{table}[H]
\centering
\scriptsize
\begin{tabular}{p{0.08\linewidth}p{0.30\linewidth}p{0.28\linewidth}p{0.10\linewidth}}
\hline
ID & Objective & Final evaluation & Status \\
\hline
O1 & Implement comparable zkVM benchmark targets & The repository contains
operation, AES, Salsa20, LowMC, CTR, and CTR+HMAC targets using a shared harness
style & Achieved \\
O2 & Measure performance under one fixed harness & Main and operations run
IDs support the report figures and tables, but mostly with one trial per target
& Partial (single-trial caveat) \\
O3 & Explain LowMC/AES behavior & The report links a \(34.85\times\) LowMC/AES
proving-time gap and \(25.59\%\) LowMC optimisation gain to VM-level cost
structure & Achieved \\
O4 & Build authenticated AES pipeline & CTR-only and CTR+HMAC target families run
with matched block sizes and receipt-backed outputs & Achieved \\
O5 & Provide scoped security and limitations analysis & The report gives a threat
model, game sketch, assumptions, limitations, and mitigation roadmap, but not a
mechanized proof & Partial (proof sketch, not full formal proof) \\
\hline
\end{tabular}
\caption{Objective-by-objective success evaluation.}
\label{tab:conclusion-objective-evaluation}
\end{table}

Overall, the project meets its implementation and analysis objectives, while the
statistical-strength and formal-proof objectives remain intentionally scoped. The
most important quantitative achievement is the measured separation between AES
and LowMC in this RISC Zero implementation. The most important engineering
achievement is the final authenticated AES target split, which makes
confidentiality-only and authenticated runs directly comparable. The most
important methodological achievement is the traceable path from objective, to
artifact, to interpreted result.

\section{Personal Reflection}

The project developed a clearer appreciation of the gap between theoretical
expectation and practical engineering constraint. Starting with the assumption
that LowMC would outperform AES in a ZK context, the measured results showed how
VM-level implementation costs can invert circuit-level intuitions. Adjusting that
expectation in response to artifact-backed evidence, rather than holding to the
prior design assumption, was the most practically valuable skill the project
reinforced.

Building the benchmark harness infrastructure required more engineering investment
than initially scoped. Reproducibility, artifact traceability, and config-driven
trial management are non-trivial in a system where proving jobs take minutes and
machine state changes between runs. Managing that infrastructure while delivering
evaluation results within the dissertation timeline required prioritisation and
explicit scope decisions --- most notably the single-trial benchmark campaign --- which
themselves became part of the reported methodology rather than hidden limitations.

From a technical perspective, constructing the security argument for the final
authenticated pipeline clarified how proof validity and cryptographic data
authenticity answer different questions, and why both layers are needed in a
deployable system. Expressing that decomposition as a structured game sketch
strengthened the intellectual contribution of the project beyond a purely
empirical benchmark report, and highlighted where a fuller formal treatment would
add value in follow-up work.
